\par\smallskip
On conserve les graphies récurrentes « bidimensionelle » et
« bidimensionelles » (sommaire, p.~2 et 26), « occurence » (p.~6, 30, 35,
47 et 53), « boite » (p.~8--9), « polygône » et « polygônes » (p.~20--21
et 26), « résidual » (p.~33 et 41), ainsi que « Légo-Teichmüller ».
Les expressions « points bases »
(p.~5 et 8), « savoir » introduisant un exemple, et le titre « Histoire
de Modèles » (p.~44--45) gardent leur forme. Les usages d'accord que
l'examen a maintenus restent à leur place.
\par
Parmi les autres formes conservées après examen figurent « astéristiques »,
« arètes », « trésaillir », « crû », « degénérescence », « celà »,
« articielle », « chapître », « répéré », « téléscopage »
et « impréssionné ». Les noms « Eilenberg-Maclane », « van Kampen » et « A. Campo » suivent
la dactylographie. Les espacements de « au delà », « c'est à dire »
et « terre à terre », ainsi que le trait de « grosso-modo », sont
conservés aux endroits vérifiés. Les abréviations, les guillemets,
les soulignements et la ponctuation restent ceux du texte vérifié.
Les traits d'union aux anciennes fins de ligne sont conservés, même
lorsque la nouvelle composition déplace la coupure. Les seules
corrections de coquilles sont celles que le maître annoté et l'apparat
justifient individuellement.
